\subsection{Resistive Current Divider}

This block is constructed from two resistors in parallel and is used
to attenuate a current signal. It is one of the simplest and most
commonly used circuit in electronics.

\subsubsection{Schematic}
\begin{circuitfig}
  % Paths, nodes and wires:
  \draw (6.5, 4.75) to[european voltage source, mirror, v=$V_{in}$, i^>=$I_{in}$, current/distance=0.6] (6.5, 6.5);
  \draw (8, 6.5) to[european resistor, l_={$G_1$}] (8, 4.75);
  \draw (8, 6.5) to[european resistor, l={$G_2$}] (8, 8.25);
  \node[ground] at (6.5, 4.75){};
  \node[ground] at (8, 4.75){};
  \node[ground, xscale=-1, yscale=-1] at (8, 8.25){};
  \draw (6.5, 6.5) -- (8, 6.5) -- (8, 6.5);
  \draw[-latex] (8.5, 6.25) -- (8.5, 4.75);
  \draw[-latex] (8.5, 6.75) -- (8.5, 8.25);
  \node[shape=rectangle, minimum width=0.715cm, minimum height=0.465cm](N1) at (8.875, 5.5){} node[anchor=center] at (N1.text){$I_{G_1}$};
  \node[shape=rectangle, xscale=-1, yscale=-1, minimum width=0.715cm, minimum height=0.465cm](N2) at (9.125, 7.5){} node[anchor=center] at (N2.text){$I_{G_2}$};
  \draw[dash pattern={on 1.6pt off 1.6pt}, -stealth] (8, 6.5) -- (9, 6.5);
  \node[shape=rectangle, minimum width=0.715cm, minimum height=0.465cm](N3) at (9.375, 6.5){} node[anchor=center] at (N3.text){$Node$};
\end{circuitfig}

\subsubsection{Transfer Function}
Applying KCL at the node between \(G_1\) and \(G_2\) we get:
\begin{align*}
  I_{in} &= I_{G_1} + I_{G_2} \\
  I_{in} &= (V_{in} - 0)G_1 + (V_{in} - 0)G_2 \\
  I_{in} &= V_{in}G_1 + V_{in}G_2 \\
  I_{in} &= V_{in}(G_1 + G_2) \\
  V_{in} &= \frac{I_{in}}{G_1 + G_2}
\end{align*}
For \(I_{G_1}\):
\begin{align*}
  I_{G_1} &= (V_{in} - 0)G_1 \\
  I_{G_1} &= V_{in}G_1 \\
  I_{G_1} &= \frac{I_{in}}{G_1 + G_2}G_1 \\
  \frac{I_{G_1}}{I_{in}} &= \frac{G_1}{G_1 + G_2} \\
  \boxed{H_1(s) = \frac{N_1(s)}{D_1(s)} = \frac{I_{G_1}}{I_{in}} = \frac{G_1}{G_1 + G_2}}
\end{align*}
For \(I_{G_2}\):
\begin{align*}
  I_{G_2} &= (V_{in} - 0)G_2 \\
  I_{G_2} &= V_{in}G_2 \\
  I_{G_2} &= \frac{I_{in}}{G_1 + G_2}G_2 \\
  \frac{I_{G_2}}{I_{in}} &= \frac{G_2}{G_1 + G_2} \\
  \boxed{H_2(s) = \frac{N_2(s)}{D_2(s)} = \frac{I_{G_2}}{I_{in}} = \frac{G_2}{G_1 + G_2}}
\end{align*}
\paragraph{Total Output:}
\begin{align*}
  I_{in} &= I_{G_1} + I_{G_2} \\
  I_{in} &= I_{in} \frac{G_1}{G_1 + G_2} + I_{in} \frac{G_2}{G_1 + G_2}
\end{align*}

\subsubsection{Analysis}
\begin{enumerate}
  \item \textbf{Resistor Ratio (\(\frac{G_1}{G_2}\)):}
    We know,
    \begin{align*}
      I_{G_1} &= (V_{in} - 0)G_1 \\
      I_{G_2} &= (V_{in} - 0)G_2
    \end{align*}
    Dividing both equations we get:
    \begin{align*}
      \frac{I_{G_1}}{I_{G_2}} = \frac{G_1}{G_2} \\
      \frac{G_1}{G_2} = \frac{I_{G_1}}{I_{G_2}}
    \end{align*}
    Assume \(I_{G_1} = I_{out}\), then our KCL equation
    becomes:
    \begin{align*}
      I_{in} &= I_{out} + I_{G_2} \\
      I_{G_2} &= I_{in} - I_{out}
    \end{align*}
    Substituting \(I_{G_2}\) in the previous equation
    we get:
    \[
      \frac{G_1}{G_2} = \frac{I_{out}}{I_{in} - I_{out}}
    \]
  \item \textbf{Impedance's:}
    \paragraph{At Input:}
      \begin{align*}
        Z_{in} &= \frac{1}{G_1 + G_2} \\
        Z_{in} &= \frac{1}{\frac{1}{R_1} + \frac{1}{R_2}} \hspace{1cm} \text{(Substituting \(G_1 = \frac{1}{R_1}\) and \(G_2 = \frac{1}{R_2}\))} \\
        Z_{in} &= \frac{1}{\frac{R_1 + R_2}{R_1 R_2}} \\
        Z_{in} &= \frac{R_1 R_2}{R_1 + R_2}
      \end{align*}
    \paragraph{At Output:}
      \begin{align*}
        Z_{out} &= \frac{1}{G_1} + \frac{1}{G_2} \\
        Z_{out} &= R_1 + R_2 \hspace{1cm} \text{(Substituting \(G_1 = \frac{1}{R_1}\) and \(G_2 = \frac{1}{R_2}\))}
      \end{align*}
  \item \textbf{Loading Effect:}
    The output of the current divider is a current signal.
    If we connect a load to the output, it will change the current flowing through
    \(G_1\) and \(G_2\). The load will form a parallel combination with
    \(G_1\) or \(G_2\), depending on where it is connected. This will
    change the effective conductance seen by the input current, thus
    altering the current division ratio. To minimize loading effects,
    the load impedance should be at least 10 times larger than the output
    impedance of the current divider.
\end{enumerate}

\subsubsection{Question}
Design a resistive voltage divider to attenuate a 5A current signal to
10mA. Establish the ratio of the resistors and calculate the input
and output impedance.
\paragraph{Solution:}
\begin{enumerate}
  \item Calculate the resistor ratio (\(\frac{G_1}{G_2}\)):
    \begin{align*}
      \frac{G_1}{G_2} &= \frac{I_{out}}{I_{in} - I_{out}} \\
      \frac{G_1}{G_2} &= \frac{0.01A}{5A - 0.01A} \\
      \frac{G_1}{G_2} &= 0.002002 \\
      \frac{R_2}{R_1} &= 0.002002 \hspace{1cm} \text{(Substituting \(G_1 = \frac{1}{R_1}\) and \(G_2 = \frac{1}{R_2}\))}
    \end{align*}
  \item Calculate the resistor values:
    Assume \(R_1 = 5k\Omega\), then:
    \begin{align*}
      R_2 &= 0.002002 \cdot R_1 \\
      R_2 &= 0.002002 \cdot 5k\Omega \\
      R_2 &\approx 10.01 \Omega
    \end{align*}
  \item Calculate the input impedance:
    \begin{align*}
      Z_{in} &= \frac{R_1 R_2}{R_1 + R_2} \\
      Z_{in} &= \frac{5k\Omega \cdot 10.01 \Omega}{5k\Omega + 10.01 \Omega} \\
      Z_{in} &\approx 9.99 \Omega
    \end{align*}
  \item Calculate the output impedance:
    \begin{align*}
      Z_{out} &= R_1 + R_2 \\
      Z_{out} &= 5k\Omega + 10.01 \Omega \\
      Z_{out} &\approx 5.01 k\Omega
    \end{align*}
\end{enumerate}

